%! Author = gabry
%! Date = 03/07/2026

\chapter{Progettazione}\label{ch:progettazione}

Successivamente alla fase di analisi dei requisiti, la fase di progettazione definisce come tali requisiti saranno
effettivamente implementati, sia a livello di struttura software che di interfaccia utente.

\section{Descrizione generale}\label{sec:descrizione-generale}

\section{Scelte architetturali}\label{sec:scelte-architetturali}

\subsection{Migrazione verso Kotlin Multiplatform}\label{subsec:migrazione-verso-kotlin-multiplatform}

Il progetto è stato inizialmente sviluppato in Kotlin nativo Android, utilizzando strumenti Java e non solo, come:

\begin{itemize}
    \item EvalEx, per la lettura dei file delle espressioni per ottenere le diagnosi
    \item Firebase, per la sincronizzazione dei dati dei pazienti e delle diagnosi
    \item SpashScreen, per l'animazione iniziale di caricamento
    \item Lingua e tema di default in base al dispositivo
\end{itemize}

Questi elementi sono stati sostituiti con altre librerie simili o sostituiti, nello specifico:

\begin{itemize}
    \item EvalEx è stato rimosso, e si è invece scritto un interprete da zero nativo Kotlin
    \item Firebase è stato sostituito con la version Gitlive della community, con l'aggiunta di diverse dipendenze per i servizi Google
    \item Lo SplashScreen animato è stato rimosso in favore di uno statico
    \item La lingua e il tema di default sono stati scritti nelle rispettive iosMain e androidMain, in quanto diverse per piattaforma
\end{itemize}

\subsection{Architettura dell'applicazione}\label{subsec:architettura-dell'applicazione}
L'architettura del software si adatta agli standard di Kotlin Multiplatform (KMP). Durante la migrazione da Kotlin a Kotlin Multiplatform, l'obiettivo primario
è stato quello di massimizzare il codice condiviso tra Android e iOS, con il minor numero di dipendenze native ove possibile, con una logica di business e di presentazione (UI) unificati.

L'applicazione segue una variante del pattern architetturale \textbf{Model-View-ViewModel (MVVM)}, costruita attorno a una singola \textit{ViewModel} centralizzata (\texttt{MelanomaViewModel}).
Questa funge da singolo punto di accesso per l'intero ciclo di vita dell'applicazione, separando l'interfaccia utente dalla logica di dominio.

Le caratteristiche architetturali principali includono:

\begin{itemize}
    \item Utilizzo di Compose Multiplaform, compatibile sia con Android che iOS in Kotlin Multiplatform, tramite
    il modulo condiviso \texttt{commonMain}.
    Le viste (View) sono \textit{stateless}, e si limitano quindi a osservare i dati e a inoltrare al ViewModel le interazioni dell'utente.
    \item Utilizzo di \texttt{StateFlow} e di \texttt{MutableStateFlow} per la gestione della reattività degli stati.
    Variabili come dati anagrafici, parametri clinici e risultati della stadiazione (\texttt{StagingResult}) vengono aggiornati reattivamente.
    Viene quindi ricomposta automaticamente l'interfaccia grafica quando i flussi di dati cambiano, emettendo un nuovo stato.
    \item Navigazione tra le schermate gestita tramite stato (State-Driven), senza librerie di routing esterne (inizialmente nella versione nativa veniva utilizzato il Navigation Component di Android), viene quindi gestita in modo manuale e dichiarativo la navigazione tra le schermate.
    Una \texttt{Sealed class Screen} definisce le possibili destinazioni navigabili dell'app (\texttt{Calculator}, \texttt{PatientList} e \texttt{PatientHistory}).
    Il ViewModel mantiene lo stato della schermata attiva, e fornisce la funzione \texttt{navigateTo(screen: Screen)}.
    Cambiando lo stato della variabile, viene quindi cambiata la schermata, eseguendo eventuale logica annessa (come il caricamento dei pazienti tramite fetch).
    \item Utilizzo di \texttt{Kotlin Coroutines} per la gestione della concorrenza e dell'accesso ai dati.
    Tutte le operazioni asincrone, come il calcolo della diagnosi basato su regole dei file JSON (e tramite \texttt{StagingEngine}), inclusa la comunicazione cloud di Firebase, sono gestite
    in modo sicuro e non bloccante, anche tramite il \texttt{viewModelScope}.
    L'interazione con il database cloud Firebase, inizialmente con Firebase nativo di Google, è stata migrata alla versione KMP della community GitLive, che permette di utilizzare funzioni CRUD
    (lettura, salvataggio e cancellazione) sui dati senza dover applicare codici diversi in base al sistema operativo.
\end{itemize}

In breve, la maggior parte del codice si trova in commonMain e utilizza codice Kotlin compatibile sia con sistemi Android che iOS, con il minimo utilizzo di funzioni proprietare tramite \texttt{expect/actual}, lasciando quindi
i moduli \texttt{androidMain} e \texttt{iosMain} il più snelli possibili (salvo inizializzazioni e configurazioni iniziali obbligatorie).

\section{Specifica dei requisiti}\label{sec:specifica-requisiti}

\subsection{Interfaccia utente}\label{subsec:interfaccia-utente}

L'interfaccia utente (UI) e l'esperienza utente (UX) sono state progettate per essere intuitive e di rapido utilizzo a livello clinico,
essendo questa applicazione uno strumento di supporto per i medici, deve essere semplice e con un carico cognitivo minimo per l'inserimento dei dati
e la lettura dei risultati delle diagnosi.

Per definire gli Screen (schermate) e definire un progetto dell'applicazione prima della scrittura del codice, è stato utilizzato \textbf{Figma}.
Il prototipo è stato costruito secondo i seguenti criteri:
\begin{itemize}
    \item Parametri suddivisi logicamente in sezioni ben definite e distinte (Parametri T, N e M), racchiuse in \texttt{Card} con stili personalizzati.
    \item Pulsanti di input standard, come \textit{switch}, \textit{chip} e \textit{checkbox} selezionabili, al posto di
    opzioni più complesse come menù a tendina.
    Alcuni switch o checkbox quando premuti mostrano ulteriori parametri opzionali od obbligatori, per l'inserimento di dati specifici per le diagnosi.
    \item Supporto sia al tema chiaro che scuro, per non affaticare la vista dell'utente in ambito ospedaliero con diverse condizioni di luce.
    Viene inoltre provvisto direttamente nella barra di navigazione un pulsante per il passaggio veloce tra i temi.
\end{itemize}

\subsection{Realizzazione del layout}\label{subsec:realizzazione-layout}

Inizialmente il layout è stato realizzato seguendo un prototipo senza cSCC Testa-Collo, ma solamente con la stadiazione del melanoma TNM.
Successivamente ha subito un'estensione per supportare il nuovo tipo di diagnosi e introdurre la compatibile con Kotlin Multiplatform.
Questo è visibile anche nei prototipi Figma (Figura~\ref{fig:figma_v1} e Figura~\ref{fig:figma_v2}):

In (Figura~\ref{fig:figma_v1}), sviluppata in Kotlin Android Nativo, il layout e l'interfaccia
erano strutturati unicamente per supportare la stadiazione del Melanoma TNM, con parametri statici e senza selettori per altri tipi di diagnosi.

\begin{figure}[htbp]
    \centering
    \includegraphics[width=\textwidth]{immagini/prototipo1/screenshots/figma}
    \caption{Fase 1 (Android Nativo) - Solo Melanoma}
    \label{fig:figma_v1}
\end{figure}

In (Figura~\ref{fig:figma_v2}), con la migrazione verso Kotlin Multiplatform, si è aggiunta
anche l'opzione di diagnosi del Carcinoma a Cellule Squamose (cSCC) Testa-Collo, di conseguenza è stata aggiunta una nuova
\texttt{card}, con due pulsanti di selezione per scegliere il tipo di diagnosi (\textit{Tipo di tumore}).
Quando premuti, ricaricano dinamicamente la sezione dei parametri con i rispettivi tipi T, N e M\@.

\begin{figure}[htbp]
    \centering
    \includegraphics[width=\textwidth]{immagini/prototipo2/screenshots/figma}
    \caption{Fase 2 (KMP) - Selezione Dinamica}
    \label{fig:figma_v2}
\end{figure}

Il progetto definito in Figma è stato successivamente tradotto in codice con \textbf{Compose Multiplatform} (nella prima versione invece con \textbf{Jetpack Compose}).
Sono stati creati e utilizzati dei componenti personalizzati come \texttt{ModernChip} (un selettore) e componenti di animazione come \texttt{AnimatedContent} per ottenere delle transizioni
fluide senza interruzioni.
Grazie a questi componenti personalizzati si è resa più modulare l'applicazione e quindi anche il suo mantenimento per eventuali aggiornamenti futuri.



\section{Database}\label{sec:database}

Utilizziamo \textbf{Cloud Firebase Firestore} per la persistenza e sincronizzazione dei dati nel cloud.
Firestore è un database NoSQL a documenti offerto nell'ecosistema Firestore di Google, tuttavia dato che questo non è
direttamente supportato su Kotlin Multiplatform, come menzionato in precedenza durante la fase di migrazione, è stato necessario
adottare la libreria SDK \textit{Gitlive Firebase} della community, che offre funzionalità simili e compatibilità con Firestore e KMP.
Essa offre dei \textit{wrapper} Kotlin sia per Android che iOS, che possono essere utilizzati direttamente nel modulo \texttt{commonMain}.

Essendo un database NoSQL a documenti, esso ha un funzionamento diverso dai tradizionali database relazionali con tabelle, righe e colonne.
Firestore mantiene i dati in una struttura gerarchica di \textit{documenti} in formato simile a JSON, raggruppati in \textit{collezioni}, che a loro volta contengono
altri documenti, con le proprie raccolte.

\subsection{Modello dati}\label{subsec:modello-dati}

Utilizziamo una struttura progettata attorno a un'unica collezione root che abbiamo denominato \texttt{patients}.
I documenti al suo interno rappresentano i singoli pazienti, identificati in modo univoco tramite UUID (\textit{Universally Unique Identifier}).
Ogni paziente ha un documento che contiene campi come lo storico (un array di diagnosi) e dati anagrafici di base.

I documenti NoSQL sono generati direttamente da codice convertendo delle \texttt{data class} Kotlin in un formato compatibile con Firestore, utilizzando
la libreria \texttt{kotlinx.serialization}, che ne permette la serializzazione e deserializzazione.
Di fatto, grazie all'annotazione \texttt{@Serializable}, il framework si occupa automaticamente della traduzione dei dati in documenti Firestore,
fortemente tipizzati garantendo la sicurezza dei tipi dei dati e riducendo codice \textit{boilerplate}.

Il modello dei dati principale della classe \texttt{Patients} è la seguente:
\begin{itemize}
    \item \texttt{id}, \texttt{name}, \texttt{surname}: Stringhe che compongo i campi per l'anagrafica e l'identificazione di ogni paziente
    \item \texttt{birthDate}: Un campo opzionale, che rappresenta la data di nascita del paziente, e che viene gestita tramite la libreria \texttt{kotlinx.datetime}
    \item\texttt{history}: Un array, simile a una lista di oggetti, di \texttt{StagingRecord}
\end{itemize}

Nonostante sarebbe stato possibile modellare lo storico dei pazienti in una sotto-collezione separata, dato che stiamo lavorando
in NoSQL si è deciso di mantenerlo come un array \texttt{embedded} all'interno dei documenti dei pazienti.
Questo permette di ottimizzare prestazioni e costi che potrebbero derivare da multiple query su diverse collezioni.
Il caricamento dei pazienti e dei relativi storici, avviene automaticamente tramite la funzione \texttt{fetchPatients()} nel ViewModel,
insieme al proprio storico clinico, in una singola operazione di lettura (read).
Questa funzione viene solitamente richiamata al momento
dell'apertura della schermata dei pazienti, durante la navigazione tramite bottom bar, in una \texttt{Kotlin Coroutine}.

\subsection{Struttura record diagnosi}\label{subsec:struttura-record-diagnosi}

Gli elementi dell'array \texttt{history}, che compongono essenzialmente un'istanza della classe \texttt{StagingRecord} serializzata,
sono composti da:

\begin{enumerate}
    \item \textbf{Metadati}: Un \texttt{UUID} univoco insieme a un \texttt{timestamp} (di tipo \texttt{Instant}) che identificano e registrano il momento
    esatto in cui l'analisi è stata effettuata
    \item \textbf{\texttt{input} MelanomaInput}: Una classe che contiene i parametri clinici inseriti dall'utente, che supporta sia la stadiazione del melanoma TNM che cSCC.
    I dati includono tipi primitivi come \texttt{breslowMm} e \texttt{linfonodiPositivi}, serializzati come stringhe sul database.
    \item \textbf{\texttt{result} StagingResult}: Il risultato di output che viene generato dal motore di calcolo (\texttt{StagingEngine}), che
    memorizza i valori TNM calcolati, il follow-up suggerito e lo stadio complessivo.
    Infatti, al posto di salvare stringhe statiche per le spiegazioni, salviamo una chiave (come \texttt{explanationKeysM}), che permette di
    assegnare testi diversi in più lingue, senza salvare nel database stringhe statiche intraducibili o tradotte in più lingue, sarà poi l'interfaccia utente
    direttamente dal dispositivo a tradurre i valori chiave nel testo equivalente.
\end{enumerate}

\subsection{Interazione con Firestore}\label{subsec:interazione-con-firestore}

Il collegamento e interazione con il database vengono gestire dal \texttt{MelanomaViewModel} tramite delle \textit{Kotlin Coroutines}
permettendo di eseguire operazione \texttt{IO} senza bloccare il thread principale dell'interfaccia grafica in modo asincrono.

Gitlive Firestore, inoltre, ha delle espressioni CRUD (Create, Read, Update e Delete) molto semplici.
Operazioni come il salvataggio e l'aggiornamento dei pazienti sono convertite tramite oggetti Kotlin, per esempio:
\begin{minted}{kotlin}
db.collection("patients").document(patientId).set(newPatient)
\end{minted}
Quando viene necessario modificare solamente parzialmente dei dati (come eliminare una singola diagnosi dello storico),
viene utilizzata l'estensione \texttt{updateFields} ottimizzata da Firestore, che permette di eseguire una modifica esclusivamente su un
elemento dell'array richiesto (nel nostro caso l'array \texttt{history}) senza rieseguire una sovrascrittura dell'intero documento,
un approccio che migliora l'efficienza nella gestione dei dati in ambiente cloud.